%%%%%%%%%%%%%%%%%%%%%%%%%%%%%%%%%%%%%%%%%%%%%%%
%
% Template per Elaborato di Laurea
% DISI - Dipartimento di Ingegneria e Scienza dell’Informazione
%
% update 2015-09-10
%
% Per la generazione corretta del 
% pdflatex nome_file.tex
% bibtex nome_file.aux
% pdflatex nome_file.tex
% pdflatex nome_file.tex
%
%%%%%%%%%%%%%%%%%%%%%%%%%%%%%%%%%%%%%%%%%%%%%%%

% formato FRONTE RETRO
\documentclass[epsfig,a4paper,11pt,titlepage,twoside,openany]{book}
\usepackage{epsfig}
\usepackage{plain}
\usepackage{setspace}
\usepackage[paperheight=29.7cm,paperwidth=21cm,outer=1.5cm,inner=2.5cm,top=2cm,bottom=2cm]{geometry} % per definizione layout
\usepackage{titlesec} % per formato custom dei titoli dei capitoli
\usepackage{amsmath}
\usepackage[boxed,vlined]{algorithm2e}


%%%%%%%%%%%%%%
% supporto lettere accentate
%
%\usepackage[latin1]{inputenc} % per Windows;
\usepackage[utf8x]{inputenc} % per Linux (richiede il pacchetto unicode);
%\usepackage[applemac]{inputenc} % per Mac.

\singlespacing

\usepackage[english]{babel}

\begin{document}

% nessuna numerazione
\pagenumbering{gobble}
\pagestyle{plain}

\thispagestyle{empty}

\begin{center}
  \begin{figure}[h!]
    \centerline{\psfig{file=img/marchio_unitrento_colore_it_202002.eps,width=0.6\textwidth}}
  \end{figure}

  \vspace{2 cm}

  \LARGE{Department of Information Engineering and Computer Science\\}

  \vspace{1 cm}
  \Large{Bachelor’s Degree in\\Computer Science\\}

  \vspace{2 cm}
  \Large\textsc{Final Dissertation\\}
  \vspace{1 cm}
  \Huge\textsc{Predictive Queue Management in Supermarkets\\}
  \Large{\it{}} % Subtitle


  \vspace{2 cm}
  \begin{tabular*}{\textwidth}{ c @{\extracolsep{\fill}} c }
    \Large{Supervisors} & \Large{Student}\\
    \Large{Alberto Montresor} & \Large{Matteo Destro}\\
    \Large{Iacopo Carreras}\\
    \Large{Andrei Tamilin}\\
  \end{tabular*}

  \vspace{2 cm}

  \Large{Academic year 2019/2020}

\end{center}



\clearpage
\pagestyle{plain} % nessuna intestazione e pie pagina con numero al centro


% inizio numerazione pagine in numeri arabi
\mainmatter

%%%%%%%%%%%%%%%%%%%%%%%%%%%%%%%%%%%%%%%%%%%%%%%%%%%%%%%%%%%%%%%%%%%%%%%%%%
%%%%%%%%%%%%%%%%%%%%%%%%%%%%%%%%%%%%%%%%%%%%%%%%%%%%%%%%%%%%%%%%%%%%%%%%%%
%% Nota
%%%%%%%%%%%%%%%%%%%%%%%%%%%%%%%%%%%%%%%%%%%%%%%%%%%%%%%%%%%%%%%%%%%%%%%%%%
%% Si ricorda che il numero massimo di facciate e' 30.
%% Nel conteggio delle facciate sono incluse 
%%   indice
%%   sommario
%%   capitoli
%% Dal conteggio delle facciate sono escluse
%%   frontespizio
%%   ringraziamenti
%%   allegati   
%%%%%%%%%%%%%%%%%%%%%%%%%%%%%%%%%%%%%%%%%%%%%%%%%%%%%%%%%%%%%%%%%%%%%%%%%%
%%%%%%%%%%%%%%%%%%%%%%%%%%%%%%%%%%%%%%%%%%%%%%%%%%%%%%%%%%%%%%%%%%%%%%%%%%

% indice
\tableofcontents
\clearpage



% gruppo per definizone di successione capitoli senza interruzione di pagina
\begingroup
% nessuna interruzione di pagina tra capitoli
% ridefinizione dei comandi di clear page
\renewcommand{\cleardoublepage}{}
\renewcommand{\clearpage}{}
% redefinizione del formato del titolo del capitolo
% da formato
%   Capitolo X
%   Titolo capitolo
% a formato
%   X   Titolo capitolo

\titleformat{\chapter}
{\normalfont\Huge\bfseries}{\thechapter}{1em}{}

\titlespacing*{\chapter}{0pt}{0.59in}{0.02in}
\titlespacing*{\section}{0pt}{0.20in}{0.02in}
\titlespacing*{\subsection}{0pt}{0.10in}{0.02in}

% sommario
\chapter*{Abstract}
\label{abstract}

\addcontentsline{toc}{chapter}{Abstract}

Long waiting queues at the checkouts are one of the major causes of customer dissatisfaction in retailing. One way to increase customers' satisfaction may be to prevent waiting queues before they appear, increasing the number of open terminals according to the predicted inflow of people at the checkouts. However, keeping a checkout open has a cost in terms of manned staff, thus an optimisation which takes in account both cost and customers satisfaction is necessary.
The main objective of this thesis is to develop a predictive model of the queueing behavior in a retail store that could be used to anticipate near-term arrival rates at the checkouts, supporting the management in a more efficient pre-positioning of the staff to avoid queues. The system should also suggest an optimal checkouts configuration, minimizing the cost while maintaining the waiting times of the customers under a maximum acceptable threshold. The main hypothesis behind this approach is that it is possible to decrease the total manned time of the staff at the checkouts without any loss of customer satisfaction, i.e. without increasing the waiting time past an “unacceptable” value.
The analyses were conducted with real-world data from the RetailerIN system, that tracks in real-time the movements of the customers in a supermarket and provides insightful information about the customers’ shopping behaviour.

Il sommario dell’elaborato consiste al massimo di 3 pagine e deve contenere le seguenti informazioni:
\begin{itemize}
  \item contesto e motivazioni
  \item breve riassunto del problema affrontato
  \item tecniche utilizzate e/o sviluppate
  \item risultati raggiunti, sottolineando il contributo personale del laureando/a
\end{itemize}

\clearpage




%%%%%%%%%%%%%%%%%%%%%%%%%%%%%%%%%%%%%%%%%%%%%%%%%%%%%%%%%%%%%%%%%%%%%%%%%%
%%%%%%%%%%%%%%%%%%%%%%%%%%%%%%%%%%%%%%%%%%%%%%%%%%%%%%%%%%%%%%%%%%%%%%%%%%
%% Nota
%%%%%%%%%%%%%%%%%%%%%%%%%%%%%%%%%%%%%%%%%%%%%%%%%%%%%%%%%%%%%%%%%%%%%%%%%%
%% Sommario e' un breve riassunto del lavoro svolto dove si descrive 
%% l’obiettivo, l’oggetto della tesi, le metodologie e 
%% le tecniche usate, i dati elaborati e la spiegazione delle conclusioni 
%% alle quali siete arrivati.
%% Il sommario dell’elaborato consiste al massimo di 3 pagine e deve contenere le seguenti informazioni: 
%%   contesto e motivazioni
%%   breve riassunto del problema affrontato
%%   tecniche utilizzate e/o sviluppate
%%   risultati raggiunti, sottolineando il contributo personale del laureando/a
%%%%%%%%%%%%%%%%%%%%%%%%%%%%%%%%%%%%%%%%%%%%%%%%%%%%%%%%%%%%%%%%%%%%%%%%%%
%%%%%%%%%%%%%%%%%%%%%%%%%%%%%%%%%%%%%%%%%%%%%%%%%%%%%%%%%%%%%%%%%%%%%%%%%%      

%%%%%%%%%%%%%%%%%%%%%%%%%%%%%%%%
% lista dei capitoli
%
% \input oppure \include
%
\chapter{Introduction}
\label{cha:introduction}

This chapter consists of three sections. In the first one, the research problem behind this thesis as well as various initial hypotheses are presented. The next section gives an overview of the RetailerIN system, that collected and provided all the data used in this research. The last one describes the general structure of the thesis.

\section{Research problem and initial hypotheses}
\label{sec:research_problem_and_initial_hypoteses}

This thesis was written as part of an internship at Thinkinside\footnote{www.thinkin.io}, a company located in Trento specialized in location intelligence for indoor spaces. Thinkinside flagship product, RetailerIN\footnote{www.retailerin.com}, is designed to help stores obtain knowledge of the customers’ behavior inside their shop. Typically those stores only have information about the purchases made, but have no visibility on what the customers actually do while shopping: how they move and respond to promotions, how different displays configuration impact sales, where are the bottlenecks in the shopping flow, what are the under-visited areas… By tracking the position of the stores’ physical assets with high-precision localization sensors and by analyzing the shoppers’ movements, RetailerIN is able to derive all these valuable information and present them in a visual dashboard for the store managers.

From this analyses, the system developed by Thinkinside is also able to suggest adjusting staffing levels based on customers volume, for example by repositioning sales associates across departments or by prioritizing the activities to perform in the store in order to maximize store operations while preserving an adequate customer support. In this context, a great impact on the customers’ shopping experience is given by waiting queues. Long queues and waiting times have significant influence over the perceived service quality and therefore over the customers’ loyalty and conversion rate. Hence, avoiding, or at least reducing, the waiting queues should eventually lead to an increase on sales. In the context of supermarkets, the only way to achieve this kind of optimization is by opening and closing checkouts according to the customers flow. Typically, this is done directly by the staff, by observing the current queue length and deciding on whether a terminal shall be opened or closed. While this approach can be quite effective if the adjustments are responsive enough to sudden changes in the flow of people arriving at the checkouts, it necessarily introduces a certain delay given by the checkouts’ opening time and by some of the customers leaving a queue and joining the new one. Moreover, it does not take in account the possible traffic spikes. For example, a new checkout get opened to reduce the current queue length, but it is left unused after a short time due to the incoming traffic returning to a normal level. These issues could be prevented by defining a predictive model for the arrival rate at the checkouts, that would allow to prevent opening additional counters if the increases in traffic levels are not persistent. The previous scenario also introduces a second factor that must be considered to achieve the aforementioned optimization: the cost in terms of manned staff at the checkouts. Naturally, the minimum possible queue length can be achieved by maximizing the service capacity, that is leaving all the available terminals opened. This however would lead to overstaffing in most of the cases, since a lower number of checkouts would be able to satisfy every new arriving customer. Moreover, the overstaffing would decrease the overall store performance and increase the total time spent in idle by the staff, since it would keep them at the counters without any actual customer to serve, preventing them from focusing on the right activities at the right time. The number of open counters should therefore be minimized in order to maximize store operations while preserving an adequate customer support.

Both presented requirements, maximizing customers satisfaction by preventing queues and maximizing store operations by reducing the number of available terminals, conflict with each other, meaning that an optimal compromise must be designed. The purpose of this research is to find such compromise that should allows to control and reduce the working time beyond the checkouts without creating long waiting queues. The problem could be approached in three different ways:
\begin{itemize}
  \item Minimizing the number of open counters, such that a predetermined maximum threshold for the queue length (or waiting time) is not exceeded.
  \item Minimizing the queue length, such that a predetermined maximum threshold for the number of open counters, and therefore for the staffing level, is not exceeded.
  \item Defining a “cost function” that quantizes both factors by assigning them a value in terms of costs. By taking in account the inflow traffic, the optimal values for the two factors can be determined such that the cost function is minimized.
\end{itemize}

While the last approach could give the best results, it is difficult to assign a value, in terms of costs, to the length of a queue or to the number of open terminals, since it may depend on other external factors and specific store policies. Considering that it is easier to manipulate the number of open counters, the first approach was therefore chosen and researched. This approach is based on two main assumptions:
\begin{itemize}
  \item Optimizing the number of open checkouts in a supermarket decreases the total staff’s idle time during the shifts at the counters.
  \item It is possible to optimize the number of open checkouts without any loss on customer satisfaction, represented by the waiting queue length.
\end{itemize}

As explained before, a predictive model for the expected queue length is necessary for this optimization to be effective. This requirement introduces another hypothesis that must hold:
\begin{itemize}
  \item It is possible to forecast the waiting queue length by analyzing the customers’ shopping behavior in the previous periods.
\end{itemize}

The next chapters explore and validate these hypotheses by the analysis of real-world data from two different supermarkets.

\section{The RetailerIN system}
\label{sec:retailerin_system}

All the data used in this research is provided by the RetailerIN system. RetailerIN is an in-store analytics and engagement solution written in Scala, able to measure and analyze in real-time and with high precision (sub-meter accuracy) how shoppers move and interact with products in a physical store. RetailerIN relies on an indoor location system based on BLE tags attached to baskets and carts and on receiver antennas mounted on the ceilings. Since only the assets are tracked, all the data collected is anonymous, not traceable to specific customers. For every asset, the RetailerIN system is able to build the path of the movements and to segment it into each different customer’s shopping session, constructing a detailed “shopping journey” report. This data is then stored and processed in real-time to extract relevant indicators and metrics on the customers’ behavior. These indicators are used to build visual analytics dashboards for store and marketing managers, helping them to optimize the placement of goods, to measure the effectiveness of store layouts and to improve staff management.

\section{Structure of the thesis}
\label{sec:structure_of_the_thesis}

Chapter \ref{cha:state_of_the_art} explores the available literature and introduces different possible approaches to the problem. Chapter \ref{cha:data_analysis} describes all the analysis conducted in order to obtain a better understanding of the data as well as identify the most appropriate solutions to achieve the research goals. In Chapter \ref{cha:solution} the final prediction model is introduced and justified. Chapter \ref{cha:results} discusses the models prediction accuracy and the choices made for the final implementation. The final chapter provides a conclusion derived from the results and points out some limitations and possible future improvements of the system.

\clearpage
\chapter{State of the art}
\label{cha:state_of_art}

This chapter discuss the available related literature. First, different approaches from similar researches are described. All this researches fall into the broader discipline of \emph{Operational Research} (OR) and mainly into two specific sub-fields: \emph{Time Series Forecasting} and \emph{Queueing Theory}. The state of the art techniques for this fields are discussed in the following two chapters.

Terminology used:
\begin{itemize}
  \item \emph{Inflow/outflow rate}: number of customers that enter/exit the supermarket in a given interval.
  \item \emph{Arrival rate}: number of customers that arrive at the checkouts in a given interval.
  \item \emph{Service rate}: maximum number of customers that can be served by each checkout in a given interval.
  \item \emph{Dwell time}: time spent by a customer in the supermarket.
  \item \emph{Basket size}: number of items bought in a shopping session.
\end{itemize}


\section{Literature review}
\label{sec:literature_review}

There are various studies that try to analyze and predict the queue length in different settings, but little was found about staff optimization.

Berman et al. (2004) \cite{berman} developed a system that manages the switching of workers between a "front room", where the checkouts are located, and a "back room", that is the shop. This system process real-time data about the count of customers either in front and back rooms, with the goal of minimizing the customer's waiting time. To accomplish this performance objective, they first defined a \emph{waiting time threshold} and required it to not be exceed. Second, they defined a minimum \emph{time-average worker complement threshold} to complete all the back room work. They used a \( M/M/\infty \) queue model to determine the performance of the system with a given number of front room workers. With this model, they were able to suggest in real-time the optimal strategy of switching workers from one room to another, minimizing the average waiting time while satisfying the thresholds constraints.

Aksu H. (2018) \cite{aksu} main goal was to optimize the idle time of the staff operating at the checkout while maintaining the waiting queue length under a predetermined threshold. He proposed a model based on the inflow of customers at the entrance, the number of customers entering the checkout area, the current waiting queue length and number of available checkouts. He divided the supermarket area into a shop and a checkout area and used a video-based system to count the customers moving in each area. He then developed a prediction model to calculate a realistic forecast of customer dwell times. The inward flow into the checkout area was predicted using the inward flow into the shop, since the inward flow into the shop appear with a delay at the checkout area, and this delay is the dwell time. To consider non-standard events, the model was adjusted with real-time measured data, to increase the prediction accuracy. Queueing theory was then used to predict the average waiting time at the checkouts. Checkouts can be opened or closed based on a minimum and maximum acceptable queue length values. To avoid closing and opening too frequently, the system was able to decide in how many cases the defined waiting queue length can be exceeded, e.g. it is acceptable to exceed the target value for a short period of time.

\section{Time series forecasting}
\label{sec:time_series_forecasting}

A \emph{time series} is a collection of observations made sequentially through time. \emph{Time series forecasting} is the use of a model to predict future values based on previous observations.

A time series can be decomposed into three main components:
\begin{itemize}
  \item \emph{Trend} \( T_t \): Linear/nonlinear, increasing/decreasing behavior of the series over time.
  \item \emph{Seasonality} \( S_t \): Repeating patterns or cycles behavior over time.
  \item \emph{Residual} \( R_t \): Variability of the observations not explainable by the model.
\end{itemize}

There are two types of decomposition: \emph{additive}, where the time series would be written as \( y_t = T_t + S_t + R_t \), and \emph{multiplicative}, where \( y_t = T_t \cdot S_t \cdot R_t \).
The additive decomposition is the most appropriate if the magnitude of the seasonal fluctuations, or the variation around the trend-cycle, does not vary with the level of the time series. When the variation in the seasonal pattern, or the variation around the trend-cycle, appears to be proportional to the level of the time series, then a multiplicative decomposition is more appropriate.

Once the main components of the time series have been determined, a predictive model can be defined. There are many different time series forecasting methods available, the most popular are presented in the next sections.



\subsection{Exponential smoothing}
\label{subsec:exponential_smoothing}
In the \emph{Simple exponential smoothing} technique, the forecasted values are based on a weighted average of the previous values, where the most recent observations are given more importance using larger weights \cite{hyndman}.

A more complex implementation is the \emph{Holt-Winters exponential smoothing} approach, which is able to decompose the time series into a level, trend and seasonal component, giving a more precise forecast. However, this method is unable to model more complex series with multiple seasonality patterns, as in a supermarket inflow rate.

Taylor’s \emph{double seasonal exponential smoothing} method \cite{taylor} was developed to forecast time series with two seasonal cycles: a short one that repeats itself many times within a longer one.


\subsection{Artificial Neural Networks}
\label{subsec:artificial_neural_networks}
\emph{Artificial Neural Networks (ANN)} allows complex nonlinear relationships between the target variable and its predictors. With time series data, lagged values of the time series can be used as inputs, in what is called an \emph{Autoregressive Neural Network (AR-NN)}. With seasonal data, it is also useful to also add the last observed values from the same season as inputs \cite{hyndman}. In general, the input can be expressed as:
\[ (y_{t-1}, y_{t-2}, ..., y_{t-p}, y_{t-m}, y_{t-2m}, ..., y_{t-Pm}) \]
We use the notation \( \text{AR-NN}(p, P, k)_m \) to indicate a neural network with \( p \) previous lagged values and \( P \) lagged values from previous seasonal cycles in the input, \( k \) neurons in the hidden layer, and a seasonal cycle of \( m \) intervals.

For forecasting two steps ahead, the result of one-step forecast can be used in the input.

\subsection{ARIMA models}
\label{subsec:arima_models}
\emph{ARIMA (AutoRegressive Integrated Moving Average)} models provide another approach to time series forecasting. Exponential smoothing and ARIMA models are the two most widely used approaches to time series forecasting, and provide complementary approaches to the problem. While exponential smoothing models are based on a description of the trend and seasonality in the data, ARIMA models aim to describe the autocorrelations in the data.\cite{hyndman}

\emph{AutoRegressive} means the model uses the relationship between an observations and some number of lagged previous observations to generate a linear regression model. \emph{Integrated} refers to the use of differencing of raw observations (e.g. subtracting an observation from an observation at the previous time step) in order to make the time series stationary. A stationary time series has mean, variance and autocorrelation constant over time. \emph{Moving Average} means that the model uses the dependency between an observation and the residual error from a moving average model applied to lagged observations.

The parameters of the ARIMA model are defined as follows:
\begin{itemize}
  \item \( p \): Autoregressive order, the number of lagged observations included in the model.
  \item \( d \): Degree of differencing, the number of times that the raw observations are differenced.
  \item \( q \): Moving average order, the size of the moving average window.
\end{itemize}

A simple version of the ARIMA model can be used to forecast only non-seasonal time series, but it can be extended to model seasonal patterns with the Seasonal ARIMA (SARIMA). There are four seasonal parameters that are not part of ARIMA that must be configured:
\begin{itemize}
  \item \( P \): Seasonal autoregressive order.
  \item \( D \): Seasonal degree of differencing.
  \item \( Q \): Seasonal moving average order.
  \item \( m \): The number of time steps for a single seasonal period.
\end{itemize}
A \( \text{SARIMA}(p,0,0)(P,0,0)_m \) model is equivalent to an \( \text{AR-NN}(p,P,0)_m \) model but with a restrictions on the parameters that ensure stationarity.
However, only a single seasonal effect can be modelled with SARIMA. To model more seasonalities the SARIMAX model can be used, since it supports exogenous variables. The seasonal effects can be captured by Fourier terms and used as exogenous variables. This gives a better approximation and allows other exogenous variables to be considered (e.g. weather, holidays, ...) to further improve the accuracy.

\subsection{TBATS models}
\label{subsec:tbats_models}
An alternative approach proposed by De Livera et al. \cite{de_livera} uses a combination of Fourier terms with an exponential smoothing model and a Box-Cox transformation, in a completely automated manner. The main advantage of a TBATS model over a SARIMAX model with Fourier terms is that the seasonality is allowed to change slowly over time, while harmonic regression terms force the seasonal patterns to repeat periodically without changing. However, one drawback of TBATS models, is that they are slow to estimate, since they will consider various alternatives and fit different models.


\section{Queueing theory}
\label{sec:queueing_theory}

\emph{Queueing theory} is the mathematical study of waiting lines, or queues. A queueing model is constructed so that queue lengths and waiting time can be predicted.

Queueing theory is the mathematical study of waiting lines, or queues. A queueing model is constructed so that queue lengths and waiting time can be predicted. There are various types of queueing models. Kendall's notation is the standard system used to describe and classify these models. In 1953, Kendall [7] proposed to describe queueing models using four factors, written as $$ A/S/c/K $$, where:
\begin{itemize}
  \item $$ A $$ denotes the time distribution between arrivals to the queue
  \item $$ S $$ denotes the service time distribution
  \item $$ c $$ denotes the number of open servers
  \item $$ K $$ denotes the maximum capacity of the queue. If not specified it is assumed $$ K = \infty $$
\end{itemize}

The best representation of a supermarket queue is probably given by the $$ M/M/c $$ model, where the $$ M $$ denote a Markovian process, meaning that the inter-arrival times and the service times of the $$ c $$ servers are exponentially distributed. The customers are served in FCFS order (First Come First Served).
The utilization factor $$ \rho = \lambda/c\mu $$ describes the proportion of total service capacity being used in the system, where:
\begin{itemize}
  \item $$ \lambda $$ is the average arrival rate
  \item $$ \mu $$ is the average service rate of a single server
  \item $$ c $$ is the number of available servers
\end{itemize}

If $$ \rho >= 1 $$, i.e. the arrival rate exceeds the total service capacity, then the queue will grow indefinitely, but if $$ \rho < 1 $$ the system is stable and the average queue length can be calculated.

The main issue with this stationary approach in the context of supermarkets is that the arrival rate and the number of available terminals cannot be expressed with probability distributions with constant mean, since they are strongly time-dependent, thus non-stationary. We can denote such time-varying model as $$ M(t)/M/c(t) $$. Moreover, it could be possible that the utilization factor exceeds 1 for a short period of time (overloading), in which the queue length increases, and after that, when $$ \rho < 1 $$, the system returns slowly to its equilibrium.

A common approach for dealing time-varying rates is the Stationary Independent Period by Period (SIPP) approach, where the analysis is conducted on small intervals. For each interval, a different $$ M/M/c $$ model with constant arrival rates and constant number of available servers is created and solved with the stationary approach. However, Green et al. (2001) [8] showed that the commonly used SIPP approach is inaccurate for parameter values corresponding to many real situations, even when the time-dependent variations are small and especially for systems which operate near the critical load. Moreover, this stationary analysis requires the arrival rate to be strictly smaller than the service rate, i.e. $$ \rho < 1 $$ must hold for every interval, while many real systems can be temporarily overloaded.

Green et al. [9] proposed an easy-to-compute approximation for determining long run average performance measures for multi-server queues with periodic arrival rates, the Pointwise Stationary Approximation (PSA). This approximation was obtained by computing the expectation of the performance measure over the time interval using the stationary formula with the arrival rate that corresponds to each point in time. They showed that this approximation is an upper bound for the expected number of customers in the queue. However again, the PSA approach can only be applied to systems where $$ \rho < 1 $$ hold, i.e. temporal overloading in some intervals is strictly forbidden [10].

Rider [11] proposed an approximation of the average queue length for a $$ M/M/1 $$ model with time-varying parameters. He found that the average queue length $$ Q(t) $$ satisfies the equation:

$$ Q(t) = \lambda(t) - \mu(t) [1 - P_0(t)] $$

where $$ P_0(t) $$ is the probability of the server being idle. This equation can be interpreted to mean that the queue will increase with the arrival rate and decrease with the service rate multiplied by an "efficiency factor" $$ 1 - P_0(t) $$, i.e. the probability that the server is busy at time $$ t $$. $$ P_0(t) $$ is not a priori known, thus an approximation must be used.

Stolletz [10] proposed an improvement of the SIPP approach, the Stationary Backlog-Carryover (SBC) approach. This approach was designed for systems with temporal overloading, thus the models of consecutive periods are no longer independent from each other. Contrary to the SIPP approach that applies the $$ M/M/c/\infty $$ model to each interval, SBC utilizes the $$ M/M/c/c $$ model, also known as the Erlang’s B model, meaning that a maximum number of  $$ c $$ customers can be in the queue at any time and any further arrivals to the queue are considered blocked (i.e. lost) and carried over into future periods. This is obtained by defining an artificial arrival rate $$ \widetilde{\lambda}(t) $$ that consists of both the original arrival rate $$ \lambda(t) $$ and a backlog rate $$ b(t) $$ of the previous period, that is the rate of blocked customers leaving the $$ M/M/c/c $$system at time $$ t $$.

Once the average queue length has been determined with one of the aforementioned methods, the average waiting time in the queue can be calculated using Little’s Law [12]. This law states that the average number of customers in a stationary system $$ L_s $$ is equal to the arrival rate $$ \lambda $$ multiplied by the average time $$ W_s $$ that a customer spends in the system:

$$ L_s = \lambda W_s $$

\clearpage
\chapter{Data analysis}
\label{cha:data_analysis}

This chapter illustrates the analyses conducted to get a better understanding of the nature of the problem and the available data. The first section describes the two datasets used in these analyses and the differences between them. The following sections present a series of considerations on the distributions and patterns found in the data and on the possible applications of these findings for reaching the objectives of this thesis.

\section{Datasets description}
\label{sec:datasets_description}

The datasets were generated by the RetailerIN system described in Section~\ref{sec:retailerin_system}. This system tracks, collects and analyses the movements of carts and baskets utilized by the customers. For every shopping session of a customer using a tracked asset, different information are collected and organized:
\begin{itemize}
  \item A unique \emph{session identifier}.
  \item An \emph{asset identifier} that can be used to trace the physical asset used. There are two types of available assets: \emph{carts} and \emph{baskets}.
  \item The \emph{start} and \emph{end} time of the session, i.e. when the customer enters and leaves the shop.
  \item The total \emph{dwell time} of the session (in seconds), i.e. the total time spent by the customer in the shop.
  \item An estimate of the total \emph{movement time} and the \emph{stationary time}, such that \( \text{dwell time} = \text{movement time} + \text{stationary time} \).
  \item The total \emph{distance} covered while shopping (in meters).
  \item The \emph{terminal number} in which the payment is made. There are two types of terminal: the classical \emph{assisted rolling belt terminals} and the \emph{self-checkouts}.
\end{itemize}
The datasets are exported in CSV format (Comma Separated Values). All the collected information are completely anonymous and not traceable to a specific customer.

Two distinct datasets were analyzed from two distinct supermarkets, that we would call \emph{supermarket A} and \emph{supermarket B}. Supermarket A is a large-sized North American retail store, while supermarket B is a medium-sized store in northern Italy. Each dataset covers different time periods and have some additional information that are not available for the other:
\begin{itemize}
  \item \emph{Dataset A} covers a six-month period, from the 1st of September 2019 to the 29th of February 2020, for a total of \( 235788 \) unique shopping sessions. Additional data about the products bought in each session was available, in particular:
        \begin{itemize}
          \item The \emph{total value} of the items bought (in dollars)
          \item The \emph{total number} of items bought, also called the \emph{basket size}
        \end{itemize}
  \item \emph{Dataset B} covers a two-month period, from the 6th of February to the 25th of March 2020, for a total of \( 32685 \) unique shopping sessions. Additional data about the length of every waiting queues and the average waiting times for each counter was available with a one-minute frequency. Specifically, for each interval, the measures available were:
        \begin{itemize}
          \item The \emph{average number of customers in queue}, that is the mean of the queue lengths in every instant.
          \item The \emph{total number of customers in queue}, that is the total count of customers that have joined the queue.
          \item The total number of \emph{completed sessions}, that is the total number of customers that completed the purchase and left the queue.
          \item The average \emph{time in queue} spent by the customers.
        \end{itemize}
\end{itemize}

The session dataset was binned in a 10-minute interval. This interval size was chosen because it gave a good compromise between accuracy and frequency: smaller intervals would have included too much noise and randomness to obtain significant results, while larger intervals would have not allowed to implement a fast response to sudden changes in traffic levels. In this chapter only the analyses conducted on the dataset A are shown, since the greater amount of available data gave more insightful results. The dataset B is used when the forecast accuracy of the predictive models is discussed in Chapter~\ref{cha:results}.

\section{Dwell time, inflow rate and service time distributions}
\label{sec:dwell_time_inflow_rate_service_time_distributions}

First, the measured inflow rates and dwell times were grouped by day and hour to get an overview of how the customers’ behavior changes during the week.

Since carts and baskets are used by different types of customers with different shopping behavior, they were considered separately. The heatmaps in Figure~\ref{fig:inflow_rate_heatmap} and~\ref{fig:dwell_time_heatmap} show the obtained results, confirming that carts and baskets are indeed used in very different ways. Moreover, it is already possible to identify an initial cyclic behavior: for example, the carts present an higher usage during weekends. These seasonal cycles are further analyzed in the next section.

\begin{figure}
  \begin{center}
    \inputpgf{img}{inflow_rate_heatmap.pgf}
  \end{center}
  \caption{The heatmap showing the distribution of the inflow rates, in customers/10min.}
  \label{fig:inflow_rate_heatmap}
\end{figure}

\begin{figure}
  \begin{center}
    \inputpgf{img}{dwell_time_heatmap.pgf}
  \end{center}
  \caption{The heatmap showing the distribution of the dwell times, in minutes.}
  \label{fig:dwell_time_heatmap}
\end{figure}

As explained in Section~\ref{sec:queueing_theory}, a lot of queueing theory models assume that the parameter values follow different specific probability distributions. The following distributions analyses were conducted to verify the applicability of these models.

First, the distribution of the dwell time shown in Figure~\ref{fig:dwell_time_distribution} was investigated and approximated with an \emph{Erlang distribution}. Again, carts and baskets were considered separately for the same reasons stated previously.

\begin{figure}
  \begin{center}
    \inputpgf{img}{dwell_time_distribution.pgf}
  \end{center}
  \caption{The probability density function of the dwell time, with the corresponding Erlang distribution fit (in red).}
  \label{fig:dwell_time_distribution}
\end{figure}

For the inflow rate, the distribution of the \emph{inter-entry time} was calculated, shown in Figure~\ref{fig:inter_arrival_time_distribution}. The inter-entry time is the time passed between the entry of two consecutive customers. In this case, the distribution can be approximated by an \emph{exponential distribution}.

\begin{figure}
  \begin{center}
    \inputpgf{img}{inter_arrival_time_distribution.pgf}
  \end{center}
  \caption{The probability density function of the inter-entry times with the corresponding exponential distribution fit (in red).}
  \label{fig:inter_arrival_time_distribution}
\end{figure}

Next, the distribution of the checkouts’ service time was investigated. Since there are no available measurements of the actual service time, the \emph{basket size} \( n \), i.e. the total number of products bought in each session, was used as an approximation. Given the service time \( S \), we can define it as combination of:
\begin{itemize}
  \item a \emph{service time per item} \( S_{item} \);
  \item an \emph{extra service time} \( S_{extra} \), to take in account possible additional time spent at the checkouts, for example for complete the payment;
\end{itemize}
and we can write:

\begin{equation}
  S = n S_{item} + S_{extra}
\end{equation}

How the approximated values for \( S_{item} \) and \( S_{extra} \) are calculated is described in the next chapters. Since we considered constants \( S_{item} \) and \( S_{extra} \), the service time distribution follows the basket size distribution, shown in Figure~\ref{fig:basket_size_distribution}.

\begin{figure}
  \begin{center}
    \inputpgf{img}{basket_size_distribution.pgf}
  \end{center}
  \caption{The probability density function of the basket size with the corresponding Erlang distribution fit (in red).}
  \label{fig:basket_size_distribution}
\end{figure}

\section{Time series analysis}
\label{sec:time_series_analysis}

Different time series analysis techniques were used to extract the main components of the time series, such as trends and seasonalities. With this decomposition it was possible to verify whether additional data preprocessing were needed and to identify specific patterns and features of the time series that could be exploited by the forecasting process.

\subsection{Inflow rate autocorrelation}
\label{subsec:inflow_rate_autocorrelation}

In order to identify the number of previous values that directly influence the current value, the \emph{autocorrelation} of the inflow rate has been calculated. This measurement is useful to identify a first set of relevant features that could be used in a predictive model.

\emph{Autocorrelation} refers to the degree of correlation between the values of the same variables across different past observations in the data. It measures the linear relationship between a variable's current value and its past values. The autocorrelation values can range from -1 to +1, where +1 represents a perfect positive linear relationship (an increase in one of the values corresponds to a proportional increase in the other), -1 represents a perfect negative linear relationship, and 0 represents the absence of any linear relationship.

When data have a trend, the autocorrelations at small lags tend to be large and positive because observations nearby in time have also similar values. Therefore, the \emph{autocorrelation function} (ACF) of a trended time series tend to have positive values that slowly decrease as the lags increase. The inflow rate and dwell time trends are investigated in the next section.

When data have a seasonality, the autocorrelations will be larger for the seasonal lags (multiples of the seasonal frequency) than for the other lags~\cite{hyndman2018}. With the analysis of the inflow rate ACF in Figure~\ref{fig:inflow_rate_acf}, it was clear that the time series present two main seasonalities:
\begin{itemize}
  \item a \emph{daily seasonality}, with a cycle length of \( {\sim}144 \) lags (\( {\sim}24 \) hours), meaning that the inflow traffic peaks are repeated every day at nearly the same hours;
  \item a \emph{weekly seasonality}, with a cycle length of \( {\sim}1008 \) lags (\( {\sim}7 \) days), meaning that each day of week has a distinct traffic pattern that is repeated every week.
\end{itemize}

\begin{figure}
  \begin{center}
    \inputpgf{img}{inflow_rate_acf.pgf}
  \end{center}
  \caption{The inflow rate's ACF with the daily and weekly seasonalities.}
  \label{fig:inflow_rate_acf}
\end{figure}

The \emph{partial autocorrelation} (PACF) is the autocorrelation between two values after removing any linear dependence on the values between them. For example, the autocorrelation at lag 1 is the coefficient of correlation between \( y_t \) and \( y_{t-1} \), which is probably also the correlation between \( y_{t-1} \) and \( y_{t-2} \). But if \( y_t \) is correlated with \( y_{t-1} \), and \( y_{t-1} \) is equally correlated with \( y_{t-2} \), then we should expect to find correlation also between \( y_t \) and \( y_{t-2} \). The partial autocorrelation excludes this mutual relation from the correlation coefficient, and, at lag \( k \), it can be described as the correlation between \( y_t \) and \( y_{t+k} \) after removing the effects of \( y_{t+1}, ..., y_{t+k-1} \). The analysis of the partial autocorrelation function, shown in Figure~\ref{fig:inflow_rate_pacf}, is useful to get an insight about the number of lagged values that should be considered by a predictive regression model. Moreover, a positive autocorrelation denotes a \emph{non-stationary} time series, that is a time series whose statistical properties change over time, i.e. the trend and seasonality affect the value of the time series at different times. More precisely, if \( \{ y_t \} \) is a stationary time series, then for all \( s \), the distribution of \( ( y_t, ..., y_{t+s}) \) does not depend on \( t \)~\cite{hyndman2018}. The majority of the forecasting methods presented in Section~\ref{sec:time_series_forecasting} are based on the assumption that the time series is stationary, so the data must be \emph{differenced} before training the model to ensure stationarity. The differencing process used is explained in details in Section~\ref{subsec:artificial_neural_network_model}.

\begin{figure}
  \begin{center}
    \inputpgf{img}{inflow_rate_pacf.pgf}
  \end{center}
  \caption{The inflow rate's PACF. The positive correlation values denote a non-stationary time series.}
  \label{fig:inflow_rate_pacf}
\end{figure}

\subsection{Time series decomposition}
\label{subsec:time_series_decomposition}

Decomposing the time series into its main components is useful to find additional patterns and to validate the previous results. While in the previous section two different seasonal cycles were determined, the main goal of this section is to identify any possible trend variation among different seasons that should be taken in account when designing the final predictive model. As described in Section~\ref{sec:time_series_forecasting}, there are two types of decomposition: \emph{additive} and \emph{multiplicative}. Since the variations around the seasonal components were not proportional to the level of the time series, the additive decomposition was chosen.
First, the \emph{trend} of the inflow rate was calculated using a moving window average with a size of \( 1008 \) intervals (that is equal to one week of observations). As shown in Figure~\ref{fig:inflow_rate_trend}, the general trend seems to be stable among different months, except for September where there seems to be a greater inflow traffic. With more data available, it should be possible to identify additional monthly or yearly seasonal cycles that could be exploited to further improve the forecasting model accuracy.

\begin{figure}
  \begin{center}
    \inputpgf{img}{inflow_rate_trend.pgf}
  \end{center}
  \caption{The inflow rate's trend.}
  \label{fig:inflow_rate_trend}
\end{figure}

The same analysis was then conducted for the dwell time (see Figure~\ref{fig:dwell_time_trend}). The dwell time’s trend is not as stable as the inflow rate’s, since a slightly increase is visible during November and December.

\begin{figure}
  \begin{center}
    \inputpgf{img}{dwell_time_trend.pgf}
  \end{center}
  \caption{The dwell time's trend.}
  \label{fig:dwell_time_trend}
\end{figure}

\subsection{Service time approximations}
\label{subsec:service_time_approximations}

As described in Section~\ref{sec:dwell_time_inflow_rate_service_time_distributions}, an approximation of the checkouts' service time \( S \) can be obtained using the basket size by:

\begin{equation}
  S = n S_{item} + S_{extra}
\end{equation}

However, this approximation is feasible only if the number of products bought in each session is available. Unfortunately, that is not the case for every store, so other approximations were investigated.

An alternative approach is to use the average number of customers served by each terminal in a time interval. Since the service rate should represent the maximum service capacity of each terminal, the periods where the counters spent some time in \emph{idle}, i.e. without serving any customer, had to be excluded from the approximation. Therefore, only the intervals with a measured queue length greater than a predetermined threshold were considered. This threshold was determined by analyzing the measured \emph{outflow rate} and comparing it to the \emph{total service rate} (with \(\text{total service rate} = \text{service rate} \cdot \text{number open terminals} \)). When the two measurements have similar values, the maximum service capacity has been reached and the checkouts have probably spent no time in idle, thus the number of customer served in that interval should give an accurate approximation of the actual service rate. However, again, this approach is implementable only if the dataset contains the queues’ status data.

To get an approximation that could be used in every setting, with data that is available independently from the store, the relation between the service and dwell times was researched. As shown in the Section~\ref{sec:dwell_time_inflow_rate_service_time_distributions}, both measurements follow an erlang distribution. Intuitively, a longer dwell time should correspond to a greater amount of products bought (i.e. a greater basket size) and therefore a longer service time. For this reason, the basket size of a session could be used to get an approximation of the expected service time. To understand this relation, the correlation between the two measurements was calculated, but the results confuted this hypothesis, since the correlation coefficient resulted in a value of approximately \( 0.24 \). Moreover, Berman et al.~\cite{berman} showed that there is no such relation between these information. This means that the dwell time has no influence over the service time, probably because the majority of the dwell time is spent looking and deciding and there could be other unrelated factors that influence the service time, such as the use of coupons and credit cards during the payment.

The last approach presented was given by the analysis of the time spent in idle by each terminal \( t_i \), that is the time in which a terminal does not have any customer to serve. To do so, the \emph{inter-exit times} \( \Delta t_e \) were calculated, that are the times passed between the consecutive exits of two customers served by the same counter:
\begin{itemize}
  \item if \( \Delta t_e > S \), the checkout has probably spent some time in idle, in particular \( t_i = \Delta t_e - S \). Moreover, if \( \Delta t_e \gg S \) the checkout is probably closed for that interval;
  \item otherwise, if \( \Delta t_e \le S \), by the time the session’s purchase has been completed by the checkout, another customer is ready to be processed. In this case there is probably a waiting queue for that counter.
\end{itemize}

The hours with peak outflow traffic, i.e. with minimum inter-exit rates, were taken in consideration to approximate the value of \( S \). That is because when the inter-exit times are low, there is probably a queue and consequently the checkouts should have spent no time in idle, meaning that \( t_i = 0 \) and \( S = \Delta t_e \). If the data about the products bought in each session is available, the approximated service time \( S \) thus obtained could be used to determine the values of \( S_{item} \) and \( S_{extra} \) introduced before.

\subsection{Service rate and open terminals count approximation}
\label{subsec:service_rate_and_open_terminals_count_approximation}

Once an approximation of the service time \( S \) has been determined, the \emph{service rate} \( \mu \), that is the number of customers served in a time interval \( \Delta t \), can be obtained as:
\begin{equation}
  \mu = \frac{\Delta t}{S}
\end{equation}

Intuitively, the total service rate should be directly influenced by the number of open terminals, considering that with more terminals available the total service capacity of the checkouts increases. For this reason, the relation between these two measurements were investigated. Since a direct measurement of the total service rate was not available and it is equal to the total number of customers that leave the checkout area and that, after a small delay, exit the store, the outflow rate was used as approximation. The correlation between the two measures resulted in a value of approximately \( 0.92 \), meaning that there is a strong linear relationship between them.

To verify this relation, the mean of the service rates per terminal \( \bar{\mu} \) was used to obtain the outflow rate \( o(t) \) as:
\begin{equation}
  o(t) = c(t) \cdot \bar{\mu} = c(t) \cdot \frac{\Delta t}{\bar{S}}
\end{equation}

A comparison between the approximated and real outflow rate can be seen in Figure~\ref{fig:outflow_approximation}. The accuracy of the results shows that a time-independent service time can be considered, since a constant value for \( S \) gives a good approximation of the actual outflow.

\begin{figure}
  \begin{center}
    \inputpgf{img}{outflow_approximation.pgf}
  \end{center}
  \caption{The real outflow rate with the corresponding service-rate-based approximation.}
  \label{fig:outflow_approximation}
\end{figure}

\clearpage
\chapter{Methodology}
\label{cha:methodology}

This chapter describes the methods used to build the prediction model. The main goal is to optimize the number of open checkouts while maintaining the waiting times under a predetermined threshold.

The final model is composed of different submodels that cooperates sequentially. First, the measured inflow rate for the recent past is combined with a forecasted inflow rate of the immediate future. These values, combined with the dwell time forecast, are used to obtain a prediction of the arrival rate at the checkouts. The arrival rate, the number of open checkouts and an approximation of the service rate are then used with queueing theory methods to calculate the expected queue length for a given interval. Different values for the number of available checkouts can be used to get the minimum staff allocation that respects the maximum queue length threshold.

\section{Inflow rate forecast}
\label{sec:inflow_rate_forecast}

The inflow rate can be expressed as a time series, hence time series forecasting models were used. As described in Chapter 2.2, different forecasting methods were analysed and the best was selected based on the performance obtained. A comparison of the results and the performance obtained is presented on Chapter 6.2.

\subsection{Persistence model}
\label{subsec:persistence_model}
The first model tested was used to define a baseline in performance. This gives an idea of how well the other models could perform on the time series, and if a model’s performance is worse than this baseline it should not be considered. The baseline was obtained by a persistence model, were the last measured value $$ y_{t-1} $$ is used as forecast:
$$ \hat{y}_t = y_{t-1} $$

\subsection{Naïve model}
\label{subsec:naive_model}
As seen in Chapter 3, the arrival rate time series presents a strong seasonal component, so the forecast should be based on this repeating pattern. The forecast value was calculated as the average of the values in the previous weeks for the same day and time:

$$ \hat{y}_t = \hat{f}(t) = \frac{1}{N} \sum_{i=1}^{N} y_{t-iW} $$

where $$ W $$ is the number of intervals in a week and $$ N $$ is the number of previous weeks considered.

To take in account the variations in the inflow rate that can be caused by external factors, e.g. holidays, a local drift was calculated, using the errors of the forecasted values for the immediate past:

$$ \hat{y}_t = \hat{f}(t) + \frac{1}{M} \sum_{i=1}^{M} y_{t-i} - \hat{f}(t-i) $$

were $$ M $$is the number of previous steps to be considered in the drift.


\subsection{Artificial Neural Network model}
\label{subsec:ann_model}
As described in Chapter 2.2.2, Artificial Neural Networks are a popular and effective time series forecasting technique. We can denote an autoregressive neural network as $$ \text{AR-NN}(p, P, k)_m $$.
Since the inflow rate values present a strong correlation with the previous values, as seen in Chapter 3.3.2, the time series is non-stationary, so the dataset must be rendered stationary before training the model. This was done by differencing the time series, i.e. by calculating the difference between consecutive observations (also called first-order differencing). Each value of the differenced time series can be written as $$ y'_t = y_t - y_{t-1} $$. This transformation can help stabilise the mean of a time series, and therefore reduce the effects of trend and seasonalities. Since the time series present a strong seasonal component, a seasonal differencing was also applied, written as $$ y'_t = y_t - y_{t-m} $$ where $$ m $$ is the number of lags in a seasonal cycle. Since the daily seasonality was taken in account, $$ m = 144 $$. Both differences were applied to obtain stationarity, so the final time series used to train the the model was:

$$ y''_t = y'_t - y'_{t-1} = (y_t - y_{t-m}) - (y_{t-1} - y_{t-m-1}) $$

Various configuration parameters were tested and the model with the best performance was used.

\subsection{ARIMA model}
\label{subsec:arima_model}

\section{Arrival rate forecast}
\label{sec:arrival_rate_forecast}
Once an inflow rate forecasting model was defined, the measured and predicted inflow rates were combined and used to get a prediction of the arrival rate at the checkouts. The inflow rate into the shop appears with a delay at the checkout area and this delay is the customer's dwell time [2]. For such a prediction, a time-dependent probability density function $$ p_t(\tau) $$ is required, where $$ \tau $$ is the expected dwell time of a customer.

In Chapter 3.1 the general distribution of the dwell time was presented and with the analysis of Chapter 3.2 it is clear that the average dwell time values are strongly time-dependent, thus it can be seen as a stochastic process $$ \{ X_t \} $$ with $$ X_t \sim \text{Erlang}(k_t, n_t) $$, where $$ k_t = \text{E}[X_t]^2 / \text{Var}[X_t] $$ and $$ n_t =  \text{E}[X_t] / \text{Var}[X_t] $$. Given the weekly and daily seasonality of the time series, different $$ X_t $$ distributions are defined for each interval on each day of the week, for a total of $$ 7 \cdot 24 \cdot 6 = 1008 $$ distributions, since intervals of 10 minutes are considered. Each aforementioned distribution was obtained by calculating the values for $$ \text{E}[X_t] $$ and $$ \text{Var}[X_t] $$ from the past observations of the same interval and day. A maximum dwell time value $$ \tau_{max} $$ has also to be defined, either by a customizable parameter or by analysing the dwell time distribution (e.g. using the 95th percentile).

Let $$ \lambda(t) $$ be the predicted arrival rate for the time interval $$ t $$, $$ p_t(\tau_i) $$ the probability of having a dwell time of $$ \tau_i $$ such that $$ (i-1) \Delta t < \tau_i \leq i \Delta t $$, $$ \Delta t $$ the intervals size and $$ y_{t-i} $$ the measured inflow rate at time $$ t-i $$. We can write:
$$ \lambda(t) = \sum_{i=1}^{K} y_{t-i} \cdot p_t(\tau_i) $$
where $$ K $$ is such that $$ \tau_K \leq \tau_{max} \leq \tau_{k+1}  $$. However, this equation is applicable only if $$ y_{t-1}, y_{t-2}, ..., y_{y-K} $$ are known, thus only when $$ t \leq t_{now} $$. If $$ t > t_{now} $$, the forecast of the inflow rate described in Chapter 4.1 must be used. For this reason, we define the function $$ \text{in}(t) $$ as:


where $$ y_t $$ and $$ \hat{y}_t $$ are respectively the measured and predicted inflow rate at time $$ t $$. The previous equation is therefore rewritten as:

$$ \lambda(t) = \sum_{i=1}^{K} \text{in}(t-i) \cdot p_t(\tau_i) $$

For a multi-step forecast of the next $$ N $$ intervals, we can write a linear system of equations:

\[
  \begin{bmatrix}\hat{\lambda}(t) \\ \hat{\lambda}(t+1) \\ \vdots \\ \hat{\lambda}(t+N)\end{bmatrix}
  =
  \begin{bmatrix}
    y_{t-1}           & y_{t-2}           & \cdots & y_{t-(N+1)} \\
    y_{t}             & y_{t-1}           & \cdots & y_{t-N}     \\
    \hat{y}_{t+1}     & y_{t}             & \cdots & y_{t-(N-1)} \\
    \vdots            & \vdots            & \ddots & \vdots      \\
    \hat{y}_{t+(N-1)} & \hat{y}_{t+(N-2)} & \cdots & y_{t}
  \end{bmatrix}
  \begin{bmatrix}p_t(\tau_1) \\ p_t(\tau_2) \\ \vdots \\ p_t(\tau_N)\end{bmatrix}
\]

\section{Queue length forecast}
\label{sec:queue_length_forecasting}
This chapter describes the implemented model for forecasting the queue length at the checkouts. As described in Chapter \ref{sec:queueing_theory}, with the M/M/c model the required variables are:
\begin{itemize}
  \item \( \lambda(t) \): Arrival rate at the checkouts.
  \item \( \mu(t) \): Service rate of the checkouts.
  \item \( c(t) \): Total count of available (open) checkouts.
\end{itemize}
A forecast of the arrival rate has been proposed in the previous chapter and an approximation of the service rate is presented in the next chapter. The total count of available checkouts is a tunable parameter that can be use to evaluate the performance of different checkouts configurations.

\subsection{Service rate approximation}
\label{subsec:service_rate_approximation}
In Chapter 3.3.3 various possible approximations of service times were introduced. In Chapter 3.3.4 it was shown that a constant, and therefore time-independent, service rate could be used without losing too much precision. This constant service rate was obtained by the analysis of the time spent in idle by each terminal. The hours with peak outflow traffic were determined by the analysis of the queues length in each interval.

\subsection{Queue length approximation}
\label{subsec:queue_length_approximation}
As described in Chapter 2.3 the standard queueing theory does not support time-varying arrival rates and temporal overloading. The only approach that seemed to fit the context of this research was the SBC (Stationary Backlog-Carryover) approach proposed by Stolletz [10]. SBC is based on a $$ M/M/c/c $$ model and uses an artificial arrival rate $$ \widetilde{\lambda}_t $$ to take in account the temporal overloading of the queue, allowing standard queueing models to be used. The $$ \widetilde{\lambda}_t $$ value consists of both the average arrival rate $$ \lambda_t $$ and the backlog rate of the previous period $$ b_{t-1} $$, that is the rate of customers leaving the system due to blocking in the former period.
Let $$ P_t(B) $$ be the steady-state probability of blocking for the $$ M/M/c/c $$model in period $$ t $$ with artificial arrival rate $$ \widetilde{\lambda}_t $$. The backlog rate $$ b_t $$ is given by

$$ b_t = \widetilde{\lambda}_t \cdot P_t(B) $$

This blocked customers are carried over into period $$ t+1$$, which results in the artificial arrival rate $$ \widetilde{\lambda}_{t+1} $$.
Starting with $$ \widetilde{\lambda}_1 = \lambda_1 $$ and $$ b_0 = 0 $$, we can define the artificial arrival rate recursively through:

$$ \widetilde{\lambda}_t = \lambda_t + b_{t-1} = \lambda_t + \widetilde{\lambda}_{t-1} \cdot P_{t-1}(B) $$

By doing this, the customers not served in period $$ t-1 $$ are evenly spread over period $$ t $$. The value for $$ P_t(B) $$ is obtained by the Erlang’s loss formula:

$$ P_t(B) = \frac{(\widetilde{\lambda}_t / \mu_t)^{c_t}}{c_t!\sum_{k=0}^{c_t} \frac{(\widetilde{\lambda}_t / \mu_t)}{k!} } $$

The expected servers utilization $$ \rho_t $$ can then be calculated as:

$$ \rho_t = \frac{\widetilde{\lambda}_t(1 - P_t(B))}{c_t\mu_t} = \frac{\widetilde{\lambda}_t - b_t}{c_t\mu_t} = \frac{\lambda_t + b_{t-1} - b_t}{c_t\mu_t} $$ (1)

Stolletz [10] proposed two different approximations for the expected number of customers in the queue, one based on a modified arrival rate (MAR) and the other based on $$ \rho_t $$ and $$ b_t $$.

The MAR approach is based on the analysis of a $$ M/M/c/\infty $$ queueing model with the same utilization value as the $$ M/M/c/c $$ model analyzed in the previous step. To do so, a modified arrival rate $$ \lambda_t^{MAR} $$ is chosen such that the model reach the approximated utilization $$ \rho_t $$, given by:

$$ \lambda_t^{MAR} = \rho_t c_t \mu_t $$

Applying (1) results in:

$$ \lambda_t^{MAR} = \lambda_t + b_{t-1} - b_t $$

The expected number of customers in the system $$ Ls_t $$ can then be obtained by resolving the specific formula for the $$ M/M/c/\infty $$ model:

$$ Ls_t^{MAR} = c_t\rho_t + \frac{\rho_t}{1-\rho_t} \pi_{c_t^+} $$
where $$ \pi_{c_t^+} $$ is the probability of having all the $$ c_t $$ servers occupied:

$$ \pi_{c_t^+} = \frac{(c_t\rho_t)^{c_t}}{c_t!(1-\rho_t)}\pi_0$$

and $$ \pi_0 $$ denotes the probability of having $$ 0 $$ customers in the system:

$$ \pi_0 = \left[ \sum_{k=0}^{c_t-1} \frac{(c_t\rho)^k}{k!} + \frac{(c_t\rho)^{c_t}}{c_t!} \frac{1}{1-\rho_t} \right]^{-1} $$

The alternative approach proposed approximates the expected queue length with the number of backlogged customers during a period (approximation A1).The backlog rate $$ b_t $$ is multiplied by the period length $$ \Delta t $$ to obtain the number of waiting customers at the end of period $$ t $$. The expected queue length $$ Lq_t^{A1} $$ is given by:

$$ Lq_t^{A1} = b_t \Delta t $$
and the total number of customers in the system $$ Ls_t^{A1} $$ by:

$$ Ls_t^{A1} = Lq_t^{A1} + c_t \rho_t $$

They showed that the expected queue length is often overestimated by the A1 approximation. A better approximation A2 was obtained by reducing the approximated $$ Lq_t^{A1} $$ value by the expected number of non-busy servers $$ c_t (1 - \rho_t) $$, such that:

$$ Lq_t^{A2} = max\{0, Lq_t^{A1} - c_t (1 - \rho_t)\} $$

and

$$ Ls_t^{A2} = Lq_t^{A2} + c_t \rho_t $$

Using Little’s Law [12], the expected waiting time in the queue $$ Wq_t $$ and in the system $$ Ws_t $$ can be then calculated as:

$$ Wq_t = Lq_t / \lambda^{MAR}_t$$
$$ Ws_t = Ls_t / \lambda^{MAR}_t $$


\section{Checkouts optimization}
\label{sec:checkouts_optimization}
In the previous section various approximations of the average queue length were discussed. In this section, these approximations are used to calculate the optimal number of checkouts that should be open to maintain an acceptable service level while reducing the manned staff’s idle time, thus reducing the total cost.

\subsection{Optimal checkouts configuration}
\label{subsec:optimal_checkouts_configuration}

\begin{algorithm}[H]
  $ c \gets c_t$\;
  \If{\FuncSty{Ls($c$)} $ > Ls_{max}$}{
    \While{$ c \leq c_{max} $ \KwSty{and} \FuncSty{Ls($c$)} $ > Ls_{max} $}{
      $ c \gets c + 1 $\;
    }
  }
  \Else{
    \While{$ c \ge 1 $ \KwSty{and} \FuncSty{Ls($c$)} $ < Ls_{max}$}{
      $ c \gets c - 1 $\;
    }
    $ c \gets c + 1 $\;
  }
\end{algorithm}


\subsection{Opening/closing fluctuations control}
\label{subsec:opening_closing_fluctuations_control}

\clearpage
\chapter{Results}
\label{cha:results}

The first section of this chapter briefly discusses the implementation of the prototypes used for testing and for performance evaluation. The next sections present the results obtained by the forecasting techniques described in the previous chapter and justify the choices made for the final model.

\section{Prototype implementation}
\label{sec:prototype_implementation}

Different prototypes for the models presented in Chapter \ref{cha:solution} were implemented in Python to be able to evaluate the prediction accuracy. To simplify the whole evaluation process, they were implemented and tested separately, to avoid having a model’s results influencing another’s. When the final architecture of the complete predictive model was defined by choosing the prototypes that obtained the best performances, a final implementation, with all the submodules working sequentially, was written. This implementation was designed to be able to observe the customers' behavior and produce forecasts in real-time, and to be updated with newly collected observations once a day. With this structure, the prototype can be easily ported to Scala and integrated into the RetailerIN platform without needing much refactoring.

\section{Inflow rate forecast}
\label{sec:inflow_rate_forecast_results}

This section discusses the results obtained by the forecasting methods for the inflow rate presented in Section \ref{sec:inflow_rate_forecast}, in particular:
\begin{itemize}
  \item Persistence model
  \item Drift model
  \item Artificial Neural Network model
  \item SARIMA model
\end{itemize}

Given the greater amount of available data, the \emph{dataset A} (see Section \ref{sec:datasets_description}) was chosen for the models' evaluation. Each model was trained on the first half (from September to November) and tested on the second half (from December to February) of the dataset.

The models’ prediction accuracy was evaluated by analyzing the \emph{one-step forecast} errors on the test set. A \emph{forecast error} is the difference between an observed value and its forecast and it can be written as \( e_t = y_t - \hat{y}_t \). Only the forecasts for the store’s opening hours (from 6:00 to 23:00) were considered on the evaluation, since the values outside these periods are equal to zero and therefore not relevant. Different measures were used to summarize the forecast errors:
\begin{itemize}
  \item \emph{Mean Absolute Error}: \( \text{MAE} = \text{mean}(| e_t |) \)
  \item \emph{Root Mean Square Error}: \( \text{RMSE} = \sqrt{\text{mean}(e_t^2)} \)
  \item \emph{Mean Absolute Percentage Error}: \( \text{MAPE} =  \text{mean}(| 100 e_t / y_t |) \)
\end{itemize}

Since the final implementation in the RetailerIN system will work in a real-time environment, the computational performances are another important aspect that must be taken in consideration when choosing the forecasting method. This metric is composed by two measurements: the \emph{forecast time} and the \emph{update time}, that respectively indicate the average times spent by the model computing the forecast for the next interval and updating the model with new collected data. The forecasting time is the most significant of the two, since new predictions are calculated with high frequency, given the intervals’ short length, while the model updates are executed once every day, preferably when the store is closed. The time spent for the initial training, when the system is started for the first time, is not significant and therefore not considered.

Table \ref{tab:inflow_rate_forecast_results} summarizes the results obtained by each model. The residuals analysis for the AR-NN model can be seen in Figure \ref{fig:inflow_rate_forecast_residuals}.

\begin{center}
  \captionof{table}{\label{tab:inflow_rate_forecast_results}Accuracy and performance results of the inflow rate predictive models.}
  \begin{tabular}{ l r r r r }
    \hline
    Model name     & MAE             & RMSE            & MAPE               & Forecast time (ms) \\
    \hline
    Persistence    & 5.0355          & 7.1982          & 49.2727\%          & 0.0001             \\
    Drift          & 2.6586          & 3.5376          & 15.1712\%          & 0.0023             \\
    \textbf{AR-NN} & \textbf{2.6027} & \textbf{3.4878} & \textbf{14.0093\%} & \textbf{2.1529}    \\
    SARIMA         & 3.1986          & 4.2620          & 19.8370\%          & 3.3107             \\
    \hline
  \end{tabular}
\end{center}

\begin{figure}
  \begin{center}
    \inputpgf{img}{inflow_rate_forecast_residuals.pgf}
  \end{center}
  \caption{Residuals analysis of the AR-NN model.}
  \label{fig:inflow_rate_forecast_residuals}
\end{figure}

\subsection{Conclusions}
\label{subsec:inflow_rate_forecast_results_conclusions}

The best performances in terms of forecast accuracy were obtained by the AR-NN model, and the analysis of the residuals shows that there is no information left out from the approximation. However, the Drift model was chosen to be implemented in the final architecture, since it obtained very similar results and the calculation of a new forecast is faster. Moreover, the implementation of the Drift model is much simpler than a neural network, thus it can be implemented directly without needing any additional library.

\section{Arrival rate forecast}
\label{sec:arrival_rate_forecast_results}

In Section \ref{sec:arrival_rate_forecast} a forecasting method for the arrival rate was described. To evaluate the accuracy of the aforementioned model’s forecast, a measurement of the actual arrival rate was necessary. Unfortunately this information was not directly available in neither of the two datasets, so an approximation based on the queue behavior was determined, in particular, based on the total number of customers in queue, available with one-minute frequency. When the total number of customers in queue increases, a new customer must have joined the queue. Therefore, the positive difference between consecutive intervals values should represent the number of new customers that have just arrived at the checkouts. The arrival rate \( \lambda_t \) can be then written as:
\begin{equation}
  \lambda_t = Lq_t^{TOT} - Lq_{t-1}^{TOT} + \theta_{t-1}
\end{equation}
where \( Lq_t^{TOT} \) is the total number of customers in queue at time \( t \) and \( \theta_{t-1} \) is the total number of customers served at time \( t-1 \). The obtained results were then aggregated in 10-minute intervals to obtain the same frequency as the forecast. With this approach, the \emph{dataset A} could not be used for the accuracy evaluation since it does not contain data regarding the total number of customers in queue. Instead, the evaluation was conducted over the smaller \emph{dataset B}.

Since this model is still part of the time series forecasting field, the same errors measures presented in the previous section were used. The accuracy for various multi-step forecasts is shown in the Table \ref{tab:arrival_rate_forecast_results}. The number of steps indicates the number of future intervals forecasted, e.g. a three-step forecasting gives prediction values for \( \lambda_{t+3} \), having the measured inflow rate data till interval \( t \).

\begin{center}
  \captionof{table}{\label{tab:arrival_rate_forecast_results}Accuracy results of the arrival rate predictive model.}
  \begin{tabular}{ l r r r }
    \hline
    Steps & MAE    & RMSE   & MAPE      \\
    \hline
    1     & 1.3220 & 1.8464 & 20.4156\% \\
    2     & 1.4769 & 1.9971 & 21.8374\% \\
    3     & 1.5559 & 2.1961 & 24.2532\% \\
    4     & 1.7508 & 2.4305 & 27.0171\% \\
    \hline
  \end{tabular}
\end{center}

\subsection{Conclusions}
\label{subsec:arrival_rate_forecast_results_conclusions}

The results show that the model’s prediction accuracy decreases as the number of steps ahead for which the forecast is calculated increases. This is probably caused by the error introduced by the store’s inflow rate forecast: for example, the three-step predictions are based on the inflow rate’s one-step and two-step predictions, as described in Section \ref{sec:arrival_rate_forecast}. Therefore, improving the inflow rate forecasting model of the previous section would also increase the overall performances of this model.

Another important factor that must be taken in consideration when evaluating the results is the period in which the data have been collected. As stated in Section \ref{sec:datasets_description}, the \emph{dataset B} used in the evaluation covers a two-month time span, from the 6th of February to the 25th of March 2020, and includes data from a supermarket in Italy. In that period the Coronavirus emergency imposed various security measure that greatly affected the grocery stores: limits on the number of customers in the shop at the same time, queue for entering the stores and minimum safe distances. Moreover, at the beginning of the pandemic a lot of people started panic-buying, so the traffic levels were unusually high.

A further improvement would be also given by including in the model the analysis of the customers’ movements during the shopping sessions. This would require to process in real-time the position of every customer currently in the shop, with the aim of identifying some recurrent shopping behavior that could give a forecast of the actual time that each customer will spend inside the store. This forecast could take in consideration the walking speed, the entry time, the pauses, the followed path, and any other significant information that would allow to assign each customer to a specific group of people with similar shopping behavior, and therefore with similar dwell times. The dwell times such determined would probably be much more accurate, since they would not be based only on the day and hour of the shopping session as the approximation used in this research, and would therefore improve the arrival rate forecast. However, these types of advanced analyses would require high-frequency data about the asset’s movements in each session and different forecasting techniques, that are out of the scope of this research.

\section{Queue length forecast}
\label{sec:queue_length_forecast_results}

The results obtained by the arrival rate forecast presented in the last section were used in this section to calculate a forecast of the waiting queue length. In Section \ref{subsec:queue_length_approximation} three approximations of the average queue length based on the SBC approach were presented:
\begin{itemize}
  \item A1
  \item A2
  \item MAR
\end{itemize}

Besides the arrival rate, these models also requires an approximation of the service rate in order to be used. With the approach presented in Section \ref{subsec:service_rate_approximation}, a realistic time-independent counters’ service rate \( \mu \) was calculated, with a value of:
\begin{equation}
  \mu = 2.12 \frac{\text{customers}}{\text{interval}} = 0.21 \frac{\text{customers}}{\text{min}}
\end{equation}

To evaluate the accuracy of each method, a measurement of each counter’s queue length was needed. The RetailerIN system is capable of inferring these information from the customers’ movement in the store. However, this inference is susceptible to accuracy issues, due to the nature of the problem, that must be taken in account when evaluating the forecasting performance. These issues includes:
\begin{itemize}
  \item Tracking of baskets: when a customer is using a basket and his turn of being served has arrived, the basket is left at the beginning of the rolling belt terminal. It is therefore impossible to determine when the customer has been served and has left the queue.
  \item Customers leaving the queue: it is possible that a customer leaves the queue to get additional products that he forgot or to join a different shorter queue. This kind of behavior is difficult to correctly track and can lead to inaccurate measurements.
  \item Sensors accuracy: if the counters or the waiting queues are too close from each other, the system may have some difficulties distinguishing every distinct queue.
\end{itemize}

As stated in the previous, the queues length data was available only for the \emph{supermarket B}. These values covered just one week of observations, from the 2th to the 8th of March, instead of two months as the sessions dataset. Again, since the training data can still be seen as a time series, the same accuracy measures from the previous section were used. The results are shown in Figure \ref{fig:queue_length_forecast} and in Table \ref{tab:queue_length_forecast_results}.

\begin{center}
  \captionof{table}{\label{tab:queue_length_forecast_results}Accuracy results of the queue length predictive model.}
  \begin{tabular}{ l r r r }
    \hline
    Approximation & MAE    & RMSE   & MAPE      \\
    \hline
    A1            & 0.8539 & 1.6681 & 28.7874\% \\
    A2            & 0.8049 & 1.6088 & 26.5083\% \\
    MAR           & 0.4919 & 0.9646 & 19.5708\% \\
    \hline
  \end{tabular}
\end{center}

\begin{figure}
  \begin{center}
    \inputpgf{img}{queue_length_forecast.pgf}
  \end{center}
  \caption{Example of the queue length prediction results with the A1, A2 and MAR methods.}
  \label{fig:queue_length_forecast}
\end{figure}

\subsection{Conclusions}
\label{subsec:queue_length_forecast_results_conclusions}

As explained before, the accuracy of the model is greatly affected by the accuracy of the queue length measurements. Moreover, the samples used for the evaluation were limited to only one week of data and the prediction accuracy may vary considerably depending on the considered period. There is also the store’s size factor that has a significant impact on the results: with a higher traffic level, the effects of random and unpredictable queueing behavior, as well as the errors in the measurements described before, are diminished. The store used for the evaluation is small-sized, therefore the traffic levels are quite low and the accuracy is thus affected.

\clearpage
\chapter{Conclusions and future works}
\label{cha:conclusions_and_future_works}

First, this chapter briefly describes the conclusions this work sought to address along with the research goals and the methods used to reach these goals. The next section discusses potential improvements and future researches based on the insights of this thesis.

\section{Conclusions}
\label{sec:conclusions}

The results presented in the previous chapter supports one of the initial hypotheses presented in Section~\ref{sec:research_problem_and_initial_hypoteses}, that claims that is possible to forecast the waiting queue length by the analysis of the customers’ shopping behavior during the previous periods. By defining a predictive model for the customers’ inflow rate into the shop and the dwell times, it is possible to forecast the arrival rates at the checkouts, that give an idea of the traffic levels that the checkouts will have to handle in the near future. Once the arrival and service rates are established,  by applying queueing theory it is possible to determine the evolution of the waiting queues, as well as the average time each customers will have to wait before being served. The store’s management can then take decisions regarding staff allocation based directly on these forecasts, or can rely on the suggestions given by the system that automatically determines the optimal staffing levels based on the future development of customers traffic. This thesis discuss a queue-length-based optimization, that minimize the staff at the counters while maintaining the queue length under a maximum threshold, but the algorithm can be easily adapted to reflect custom internal store’s queue management policies and staff scheduling constraints.

Considering the randomness of the customers’ behavior, the forecasts obtained with this approach are accurate enough to represent a useful tool to support decision making and resource planning. Given that RetailerIN is focused on layout and staffing optimization to increase store’s profitability, this model provide an additional value to the predictive capacity of the Thinkinside’s system and it is easily implementable, since it was designed to cooperate with the already-existing components and data formats. Moreover, the analyses conducted on this thesis are generic enough to be adapted to other contexts besides supermarkets, for example airports and banks, but also call-centers, as long as measurements for the traffic and service levels are available and the setting is characterized by at least one waiting queue with adjustable service capacity.

\section{Future works}
\label{sec:future_works}

As stated previously, the main objectives of this thesis were to develop a predictive model for how the waiting queues will evolve in the near future and to support management in adjusting staffing levels in order to maximize store operations while maintaining an adequate service level. While the accuracy of the queue length forecast could be evaluated, as described in Chapter~\ref{cha:results}, it was not possible to verify the effectiveness of the suggested checkouts optimizations due to time constraints. This could be done by implementing the final model into RetailerIN and by testing it with different supermarkets layout. By ensuring that the suggestions provided are followed by the staff and by collecting the performance metrics over a time span of different months, it would be possible to compare the performance of the influenced and uninfluenced system by analyzing the average queues lengths and amount of time the staff spent at the counters. However, this is impractical for several reasons, the first of which is that it is impossible to guarantee that the checkouts configuration suggestions will be always followed, and therefore it is difficult to obtain performance metrics that are  truly representative of the effectiveness of the model. A less significative but simpler evaluation method, without the aforementioned issue, would be to implement a simulation by using realistic arrival and service rates based on the approximations described in this thesis.

Another possible improvement would be to refine the dwell times distributions by excluding the time spent in queue from the dwell times of each session. Moreover, as already described in Section~\ref{subsec:arrival_rate_forecast_results_conclusions}, the arrival rate forecast, and therefore the overall accuracy, could be further improved by implementing a more advanced dwell time prediction by considering in real time the movements of every customers in the store.

\clearpage





\endgroup


% bibliografia in formato bibtex
%
% aggiunta del capitolo nell'indice
\addcontentsline{toc}{chapter}{References}
% stile con ordinamento alfabetico in funzione degli autori
\bibliographystyle{plain}
\bibliography{biblio}
%%%%%%%%%%%%%%%%%%%%%%%%%%%%%%%%%%%%%%%%%%%%%%%%%%%%%%%%%%%%%%%%%%%%%%%%%%
%%%%%%%%%%%%%%%%%%%%%%%%%%%%%%%%%%%%%%%%%%%%%%%%%%%%%%%%%%%%%%%%%%%%%%%%%%
%% Nota
%%%%%%%%%%%%%%%%%%%%%%%%%%%%%%%%%%%%%%%%%%%%%%%%%%%%%%%%%%%%%%%%%%%%%%%%%%
%% Nella bibliografia devono essere riportati tutte le fonti consultate 
%% per lo svolgimento della tesi. La bibliografia deve essere redatta 
%% in ordine alfabetico sul cognome del primo autore. 
%% 
%% La forma della citazione bibliografica va inserita secondo la fonte utilizzata:
%% 
%% LIBRI
%% Cognome e iniziale del nome autore/autori, la data di edizione, titolo, casa editrice, eventuale numero dell’edizione. 
%% 
%% ARTICOLI DI RIVISTA
%% Cognome e iniziale del nome autore/autori, titolo articolo, titolo rivista, volume, numero, numero di pagine.
%% 
%% ARTICOLI DI CONFERENZA
%% Cognome e iniziale del nome autore/autori (anno), titolo articolo, titolo conferenza, luogo della conferenza (città e paese), date della conferenza, numero di pagine. 
%% 
%% SITOGRAFIA
%% La sitografia contiene un elenco di indirizzi Web consultati e disposti in ordine alfabetico. 
%% E’ necessario:
%%   Copiare la URL (l’indirizzo web) specifica della pagina consultata
%%   Se disponibile, indicare il cognome e nome dell’autore, il titolo ed eventuale sottotitolo del testo
%%   Se disponibile, inserire la data di ultima consultazione della risorsa (gg/mm/aaaa).    
%%%%%%%%%%%%%%%%%%%%%%%%%%%%%%%%%%%%%%%%%%%%%%%%%%%%%%%%%%%%%%%%%%%%%%%%%%
%%%%%%%%%%%%%%%%%%%%%%%%%%%%%%%%%%%%%%%%%%%%%%%%%%%%%%%%%%%%%%%%%%%%%%%%%%

\end{document}
